%% Need to keep this so stuff works
\DocumentMetadata{tagging=on,
    pdfversion = 2.0,
    tagging-setup={math/setup=mathml-SE},
    pdfstandard=ua-2,
    lang=en-US
}

\documentclass[12pt, oneside, letterpaper]{book}
%%
%
%%
%%% DOCUMENT START
\input{preamble} % Moved to another file for ease of updating
\begin{document}
%%
%
%%
%%% Document Setup Settings
%%
\pagestyle{fancy} % Set the page style to "fancy"...
\renewcommand{\headrulewidth}{0pt}
\fancyhead{}
\fancyhead[R]{\thepage}
\fancyfoot{} % clear all footer fields
\renewcommand{\footrulewidth}{0.1pt}
\fancyfoot[L]{Truman Science Center}
\fancyfoot[C]{Key: Calorimetry/Heat Capacity}
\fancyfoot[R]{15 April 2026}
%%
%
%% Title Set up
 \styleTitle{ Key:Calorimetry/Heat Capacity}
%%
%
%%
%%% START TO DOCUMENT TEXT
%%
%
\section{The Process}

The first step is always to determine which category the compound fits into because each category has different naming conventions. 

\begin{enumerate}
    \item asdflkajsd
    \item asldkfja
    \item asdflaksjd
\end{enumerate}

\subsection{ Covalent Bonding: nonmetal bound to a nonmetal}

The nonmetals are elements found in columns 13-18 (also sometimes called 

\newpage

\section{ Convert the chemical names below to chemical formulas.} 


\flushleft{\textmd{a) Carbon Monoxide}}
\vspace{20mm}
\end{document}